\scanpage{65}
\mspage{123}
\begin{flushright}number 371\end{flushright}

$(i_h)_{1\handwritten{\leq}h\handwritten{\leq}p}$ is the strictly increasing sequence of the elements of $H$.

PROPOSITION 6.- \uline{Let} $K$ (resp. $N$) \uline{be an} $A$-\uline{module having a basis} $(e_i)_{1\handwritten{\leq}i\handwritten{\leq}n}$ (resp. $(f_j)_{1\handwritten{\leq}j\handwritten{\leq}m}$) \uline{of} $n$ (resp. $m$) \uline{elements} ; \uline{let} (for $p\handwritten{\leq}\inf(m,n)$) $(e_K)$ (resp. $(f_H)$) \uline{be the corresponding basis of} $\overset{p}{\handwritten{\bigwedge}}M$ (resp. $\overset{p}{\handwritten{\bigwedge}}N$). \uline{Let} $u:M\handwritten{\to}N$ \uline{be an} $A$-\uline{linear map}, \uline{and let us denote by} $\underline{X}$ \uline{the matrix of} $u$ \uline{with respect to the bases} $(e_i)$ \uline{and} $(f_j)$. \uline{Then the matrix of} $\overset{p}{\handwritten{\bigwedge}}u$ \uline{with respect to the bases} $(e_K)$ \uline{and} $(f_H)$ \uline{is the matrix} $(X_{H,K})$ \uline{with} $\binom{m}{p}$ \uline{rows and} $\binom{n}{p}$ \uline{columns formed from the minors of order} $p$ \uline{of} $\underline{X}$.

Indeed let $(j_k)_{1\handwritten{\leq}k\handwritten{\leq}p}$ be the strictly increasing sequence of the elements of $K\handwritten{\subset}\unsure{\handwritten{[}1,n\handwritten{]}|\handwritten{[\![}1,n\handwritten{]\!]}}$ ; we have by definition
\[
\overset{p}{\handwritten{\bigwedge}}u(e_K)=u(e_{j_1})\handwritten{\wedge}u(e_{j_2})\handwritten{\wedge}\ldots\handwritten{\wedge}u(e_{j_p})\ ;
\tag*{(10)}
\]
thus the entry of the matrix of $\overset{p}{\handwritten{\bigwedge}}u$ which is in the row with index $H$ and the column with index $K$ is the component with index $H$ of the decomposable $p$-vector $u(e_{j_1})\handwritten{\wedge}\ldots\handwritten{\wedge}u(e_{j_p})$ ; it is consequently equal to $X_{H,K}$ by virtue of proposition 5.

One says that the matrix $(X_{H,K})$ is the $p$-\uline{th} \uline{exterior power} of $\underline{X}$ and one denotes it by $\overset{p}{\handwritten{\bigwedge}}\underline{X}$.

4. \uline{Expansions of a determinant}.

Let us return to formula (4) of number 1 giving the determinant of a square matrix $\underline{X}=(\handwritten{\xi}_{ij})$ of order $n$ ; for every subset $H$ of the set $\handwritten{[}1,n\handwritten{]}$, let us denote by $H'$ the complement of $H$ in $\handwritten{[}1,n\handwritten{]}$ ; if $(i_k)_{1\handwritten{\leq}i\handwritten{\leq}p}$ (resp. $(j_h)_{1\handwritten{\leq}h\handwritten{\leq}n-p}$) is the strictly increasing sequence of the elements of $H$ (resp. $H'$), we can write (\S\ 3, number 2, lemma 3)
\[
x_1\handwritten{\wedge}x_2\handwritten{\wedge}\ldots\handwritten{\wedge}x_n
=\handwritten{\rho}_{H,H'}(x_{i_1}\handwritten{\wedge}x_{i_2}\handwritten{\wedge}\ldots\handwritten{\wedge}x_{i_p})\handwritten{\wedge}(x_{j_1}\handwritten{\wedge}x_{j_2}\handwritten{\wedge}\ldots\handwritten{\wedge}x_{j_{n-p}})
\tag*{(11)}
\]
where $\handwritten{\rho}_{H,H'}$ has the meaning defined in \S\ 5, number 3, formula (12). With the notation of number 3, we have
% archive sheet number: 64
